% ============================================================
% CHAPTER 5: Sprint 2 — Policy Engine and 100 Security Rules
% (Following Tarek's sprint chapter pattern)
% ============================================================

\chapter{Sprint 2: Policy Engine and 100 Security Rules}

\section{Introduction}

This chapter covers the second sprint (February 17--28, 2026), during which we expanded the service from 12 initial policies to the full set of 100 security rules organized across 14 categories. This sprint also introduced the configurable severity threshold system and the diff-only scanning feature. By the end of this iteration, the policy engine was capable of performing comprehensive infrastructure security analysis across AWS, Docker, and Kubernetes environments.

% ============================================================
\section{Sprint 2 Objectives}
% ============================================================

The objectives for Sprint 2 were:

\begin{enumerate}
    \item Implement the 88 additional security policies (from 12 to 100 total).
    \item Organize policies into 14 logical categories covering distinct security domains.
    \item Implement the configurable severity threshold system (CRITICAL, HIGH, MEDIUM, LOW).
    \item Add the diff-only scanning feature to filter violations to changed lines only.
    \item Ensure every violation includes a detailed message and a remediation recommendation.
\end{enumerate}

% ============================================================
\section{Sprint 2 Backlog}
% ============================================================

\begin{table}[H]
\centering
\caption{Sprint 2 — Product Backlog}
\label{tab:sprint2-backlog}
\begin{tabular}{|p{0.8cm}|p{8.5cm}|c|c|}
\hline
\textbf{ID} & \textbf{User Story} & \textbf{Priority} & \textbf{Status} \\
\hline
US-07 & As a security engineer, I can enforce 100 security policies covering AWS, Docker, Kubernetes, and Terraform. & High & Done \\
\hline
US-08 & As a security engineer, I can configure the minimum severity level at which deployments are blocked. & High & Done \\
\hline
US-09 & As a developer, I receive detailed violation messages explaining the risk and how to fix it. & High & Done \\
\hline
US-10 & As a DevOps engineer, I can scan only changed lines using Git diff to reduce noise. & Medium & Done \\
\hline
US-11 & As a security engineer, I can enable or disable specific policies through configuration. & Medium & Done \\
\hline
\end{tabular}
\end{table}

% ============================================================
\section{Use Case Refinement}
% ============================================================

\subsection{Refined Use Case: Configure Severity Threshold}

% TODO: INSERT USE CASE DIAGRAM FOR SPRINT 2
% \begin{figure}[H]
%     \centering
%     \includegraphics[width=0.85\textwidth]{figures/sprint2-usecase.png}
%     \caption{Sprint 2 — Refined use case: threshold configuration}
%     \label{fig:sprint2-usecase}
% \end{figure}

\textit{[Figure: Sprint 2 Use Case Diagram — create in draw.io]}

\vspace{0.3cm}

\begin{table}[H]
\centering
\caption{Textual description — Configure severity threshold}
\label{tab:sprint2-usecase-text}
\begin{tabular}{|p{3.5cm}|p{9.5cm}|}
\hline
\textbf{Use Case} & Configure Severity Threshold for ALLOW/BLOCK Decision \\
\hline
\textbf{Actor} & DevOps Engineer / Security Engineer \\
\hline
\textbf{Precondition} & The service is running. The actor knows the available severity levels. \\
\hline
\textbf{Postcondition} & Only violations at or above the configured severity level trigger a BLOCK decision. \\
\hline
\textbf{Nominal Scenario} &
\begin{enumerate}[nosep]
    \item The actor sends a \texttt{POST /analyze} request with \texttt{block\_threshold} set to the desired level (e.g., HIGH).
    \item The policy engine evaluates all 100 policies against the submitted code.
    \item The engine identifies all violations and their severity levels.
    \item The threshold filter retains only violations at or above the configured level.
    \item If blocking violations remain, the decision is BLOCK; otherwise, ALLOW.
    \item The response includes \textit{all} violations (for transparency) but the decision is based only on those meeting the threshold.
\end{enumerate} \\
\hline
\textbf{Alternative Scenario} &
\begin{itemize}[nosep]
    \item If threshold is set to CRITICAL, only CRITICAL violations block the pipeline. HIGH, MEDIUM, and LOW violations are reported but do not prevent deployment.
    \item If threshold is set to LOW (strictest), all violations of any severity block deployment.
\end{itemize} \\
\hline
\end{tabular}
\end{table}

% ============================================================
\section{Sequence Diagram}
% ============================================================

% TODO: INSERT SEQUENCE DIAGRAM FOR SPRINT 2
% Focus on the threshold filtering mechanism:
% Client → API → PolicyEngine.analyze() → [run 100 policies] → filter by threshold → decision
% \begin{figure}[H]
%     \centering
%     \includegraphics[width=\textwidth]{figures/sprint2-sequence.png}
%     \caption{Sequence diagram — severity threshold filtering}
%     \label{fig:sprint2-sequence}
% \end{figure}

\textit{[Figure: Sprint 2 Sequence Diagram — create in draw.io]}

% ============================================================
\section{Realization}
% ============================================================

\subsection{Policy Categories Architecture}

The 100 security policies are organized into 14 categories, each targeting a specific security domain. This organization ensures comprehensive coverage and makes the codebase navigable:

\begin{table}[H]
\centering
\caption{Complete policy organization — 14 categories, 100 policies}
\label{tab:policy-categories-full}
\small
\begin{tabular}{|c|l|c|p{5.5cm}|}
\hline
\textbf{\#} & \textbf{Category} & \textbf{Count} & \textbf{Key Checks} \\
\hline
1 & Core Policies & 12 & Public S3, hardcoded secrets, public DB, root user \\
\hline
2 & AWS S3 Advanced & 5 & Versioning, MFA delete, lifecycle, block public ACL \\
\hline
3 & AWS EC2 / Compute & 7 & IMDSv2, public IP, EBS encryption, monitoring \\
\hline
4 & AWS RDS / Database & 7 & Multi-AZ, backup retention, deletion protection \\
\hline
5 & AWS IAM & 7 & Admin policy, wildcard permissions, MFA, key rotation \\
\hline
6 & Networking / VPC & 6 & Flow logs, SSH/RDP open, default VPC \\
\hline
7 & Docker Advanced & 8 & Latest tag, healthcheck, secrets in ENV, multi-stage \\
\hline
8 & Kubernetes Advanced & 12 & ReadOnly root FS, privilege escalation, host PID \\
\hline
9 & CloudTrail / Monitoring & 8 & CloudTrail, GuardDuty, Security Hub, Config \\
\hline
10 & ELB / WAF / Route53 & 7 & HTTP listener, WAF, TLS policy, DNSSEC \\
\hline
11 & Terraform General & 5 & Remote backend, provider version, state locking \\
\hline
12 & EKS / Lambda / ECS & 7 & EKS public endpoint, Lambda secrets, ECS logging \\
\hline
13 & Messaging / Cache & 5 & SNS encryption, SQS encryption, DynamoDB \\
\hline
14 & ECR / Secrets / CloudFront & 4 & ECR scanning, Secrets Manager rotation \\
\hline
\hline
& \textbf{Total} & \textbf{100} & \\
\hline
\end{tabular}
\end{table}

\subsection{Severity Level System}

Every policy is assigned a severity level reflecting the potential impact of the vulnerability:

\begin{table}[H]
\centering
\caption{Severity levels and their definitions}
\label{tab:severity-levels}
\begin{tabular}{|c|p{3cm}|c|p{6cm}|}
\hline
\textbf{Level} & \textbf{Color} & \textbf{Policies} & \textbf{Definition} \\
\hline
\cellcolor{red!20}CRITICAL & \cellcolor{red!20}Red & 8 & Immediate exploitation risk — data exposure, credential leak \\
\hline
\cellcolor{orange!20}HIGH & \cellcolor{orange!20}Orange & 35 & Significant vulnerability — public access, missing encryption \\
\hline
\cellcolor{yellow!20}MEDIUM & \cellcolor{yellow!20}Yellow & 37 & Moderate risk — missing monitoring, non-production config \\
\hline
\cellcolor{blue!10}LOW & \cellcolor{blue!10}Blue & 14 & Best practice recommendation — missing tags, documentation \\
\hline
\cellcolor{gray!10}INFO & \cellcolor{gray!10}Gray & 6 & Informational — cost optimization, code quality \\
\hline
\end{tabular}
\end{table}

\subsection{Configurable Threshold Implementation}

The threshold system allows administrators to control how strictly the pipeline is enforced. The implementation uses a severity ordering dictionary:

\begin{lstlisting}[language=Python, caption={Severity threshold filtering in PolicyEngine}, label={lst:threshold}]
SEV_ORDER = {'CRITICAL': 4, 'HIGH': 3, 'MEDIUM': 2, 'LOW': 1, 'INFO': 0}

# Inside PolicyEngine.analyze():
threshold_level = SEV_ORDER.get(block_threshold.upper(), 1)
blocking_violations = [
    v for v in all_violations
    if SEV_ORDER.get(v.severity.value, 0) >= threshold_level
]
allow = len(blocking_violations) == 0
\end{lstlisting}

\noindent This design means:
\begin{itemize}
    \item \textbf{threshold=LOW}: All violations trigger BLOCK (strictest mode).
    \item \textbf{threshold=HIGH}: Only HIGH and CRITICAL violations trigger BLOCK. MEDIUM and LOW are reported but do not block.
    \item \textbf{threshold=CRITICAL}: Only CRITICAL violations block deployment. All others pass.
\end{itemize}

\subsection{Diff-Only Scanning}

The diff-only feature filters violations to only the lines that changed in a Git commit. This reduces noise for developers working on large codebases:

\begin{lstlisting}[language=Python, caption={Diff parsing for line-level filtering}, label={lst:diff-parse}]
def parse_diff_lines(diff: str) -> set:
    """Extract line numbers added/changed in a unified diff."""
    changed_lines = set()
    current_line = 0
    for line in diff.splitlines():
        if line.startswith('@@'):
            match = re.search(r'\+(\d+)', line)
            if match:
                current_line = int(match.group(1)) - 1
        elif line.startswith('+') and not line.startswith('+++'):
            current_line += 1
            changed_lines.add(current_line)
        elif not line.startswith('-'):
            current_line += 1
    return changed_lines
\end{lstlisting}

\subsection{Educational Violation Messages}

A distinguishing feature of our service is that every violation includes three components: \textit{what} is wrong, \textit{why} it is dangerous, and \textit{how} to fix it. For example:

\begin{lstlisting}[language=Python, caption={Example violation with educational message}, label={lst:educational}]
PolicyViolation(
    rule_id="EC2_IMDSV1_ENABLED",
    severity=PolicySeverity.HIGH,
    message="EC2 instance 'web-server' allows IMDSv1 (SSRF risk)",
    line_number=15,
    recommendation="Set metadata_options { http_tokens = 'required' } "
                   "to enforce IMDSv2"
)
\end{lstlisting}

\noindent This transforms security enforcement from a ``red flag'' experience into a \textbf{learning opportunity}: developers not only learn that their code is insecure, but they understand \textit{why} and receive a concrete fix.

% ============================================================
\section{Technologies Used}
% ============================================================

Sprint 2 relied on the same technology stack established in Sprint 1, with the addition of:

\vspace{0.5cm}
\begin{tabular}{| m{2.5cm} | m{10cm} |}
\hline

\textbf{re (stdlib)} &
Python's built-in regular expression module was used extensively in the \texttt{HardcodedSecretsPolicy} and other credential-detection policies to match patterns like AWS access keys (\texttt{AKIA[0-9A-Z]\{16\}}).
\\
\hline

\textbf{dataclasses} &
Python's \texttt{@dataclass} decorator provided lightweight data containers for the \texttt{Resource}, \texttt{NormalizedCode}, \texttt{PolicyDecision}, and \texttt{PolicyViolation} models, reducing boilerplate code.
\\
\hline

\textbf{enum (stdlib)} &
Python's \texttt{Enum} class defined the \texttt{PolicySeverity} levels (CRITICAL, HIGH, MEDIUM, LOW, INFO), ensuring type safety and preventing invalid severity values at compile time.
\\
\hline

\end{tabular}

% ============================================================
\section{Sprint 2 Interfaces}
% ============================================================

\subsection{Full Scan with Threshold Configuration}

% TODO: INSERT SCREENSHOT showing a scan with block_threshold=HIGH
% \begin{figure}[H]
%     \centering
%     \includegraphics[width=0.9\textwidth]{figures/sprint2-threshold-scan.png}
%     \caption{Sprint 2 — API response with threshold=HIGH (MEDIUM violations reported but not blocking)}
%     \label{fig:sprint2-threshold}
% \end{figure}

\textit{[Screenshot: API response showing threshold filtering]}

\subsection{100 Policies Loaded}

% TODO: INSERT SCREENSHOT showing the startup log confirming 100 policies loaded
% \begin{figure}[H]
%     \centering
%     \includegraphics[width=0.9\textwidth]{figures/sprint2-100policies.png}
%     \caption{Sprint 2 — Service startup confirming 100 policies loaded}
%     \label{fig:sprint2-100policies}
% \end{figure}

\textit{[Screenshot: Terminal showing 100 policies loaded]}

% ============================================================
\section{Conclusion}
% ============================================================

Sprint 2 transformed the service from a prototype with 12 policies into a comprehensive security scanning engine with 100 policies across 14 categories. The configurable severity threshold system gives DevOps teams granular control over pipeline behavior, while the diff-only scanning feature reduces noise for incremental changes. Every violation includes an educational message with a remediation recommendation, making the service not just an enforcement tool but a learning platform.

The codebase grew to approximately 3,000 lines of policy code in \texttt{security\_policies.py}, yet remained maintainable thanks to the Template Method pattern: each policy is a self-contained class with a single \texttt{check()} method. Adding a new policy requires only creating a new class and adding it to the \texttt{get\_all\_policies()} registry.

The next chapter will focus on building the dashboard, HTML reports, and scan history system that make the service's output accessible and actionable.
